\section{Objetivos}\label{sec:objetivos_informe_final}

% -------------------------------------------------------
% HITO 1: Planificación, motivación y marco teórico
% -------------------------------------------------------

\subsection*{Hito 1 --- Objetivo general}

Establecer las bases conceptuales, tecnológicas y de planificación que orientarán el desarrollo del sistema de gestión de préstamos e inventario del laboratorio LabDIC.

\subsection*{Hito 1 --- Objetivos específicos}

\begin{itemize}
	\item Identificar las tecnologías de base de datos, backend y frontend adecuadas para el proyecto.
	\item Analizar la motivación del proyecto y las necesidades que busca cubrir.
\end{itemize}

% -------------------------------------------------------
% HITO 2: Diseño del modelo de base de datos
% -------------------------------------------------------

\subsection*{Hito 2 --- Objetivo general}

Diseñar e implementar un modelo de base de datos relacional que soporte de forma 
consistente, segura y escalable la gestión de usuarios, catálogo de referencia, 
productos, dispositivos y préstamos del sistema LabDIC, permitiendo la operación 
concurrente de múltiples administradores y clientes, así como el registro del ciclo 
de vida de los préstamos.

\subsection*{Hito 2 --- Objetivos específicos}

\begin{itemize}
	\item Registrar la información de los usuarios del sistema, sus credenciales, roles 
	y permisos.
	
	\item Implementar el módulo de autenticación a través de campos específicos del 
	modelo de usuario y su relación con roles, sin requerir tablas adicionales.
	
	\item Centralizar los datos de referencia del sistema: roles, marcas, modelos, 
	categorías, estados y ubicaciones, en un módulo de catálogo reutilizable por los 
	demás módulos.
	
	\item Administrar las unidades físicas individuales del inventario mediante el módulo 
	de dispositivos, registrando su estado actual y el historial de cambios de estado.
	
	\item Garantizar la referencialidad entre módulos mediante el uso adecuado de 
	claves primarias y foráneas.
	
	\item Facilitar la evolución futura del esquema, apoyado en el control de versiones con Alembic.
\end{itemize}

% -------------------------------------------------------
% HITO 5: Arquitectura del sistema completo
% -------------------------------------------------------

\subsection*{Hito 5 --- Objetivo general}

Documentar la arquitectura técnica del sistema, describiendo los patrones de diseño, 
convenciones de organización y mecanismos de comunicación empleados tanto en el 
backend como en el frontend.

\subsection*{Hito 5 --- Objetivos específicos}

\begin{itemize}
	\item Describir el rol y las responsabilidades de los DTOs, repositorios y 
	controladores en el backend, y explicar la delegación de tareas entre estas 
	tres capas.
	
	\item Describir el funcionamiento del router del backend y Vue Router en el frontend, y sus integraciones con la 
	aplicación Litestar (backend) y la implementación de 
	guards de navegación basados en autenticación y roles (frontend).
	
	\item Describir el uso de Pinia para la gestión del estado 
	global persistente.
	
	\item Presentar el flujo de consumo de la API del sistema.
\end{itemize}

% -------------------------------------------------------
% HITO 3: Implementación del backend
% -------------------------------------------------------

\subsection*{Hito 3 --- Objetivo general}

Implementar el backend completo del sistema, desarrollando los servicios, repositorios, 
DTOs y controladores de los seis módulos funcionales (Usuarios, Autenticación, Catálogo, 
Productos, Dispositivos y Préstamos), exponiendo la API REST documentada 
que sirva como base para el desarrollo del frontend en los hitos siguientes.

\subsection*{Hito 3 --- Objetivos específicos}

\begin{itemize}
	\item Documentar el uso de Alembic como herramienta de migración iterativa durante 
	el desarrollo, describiendo los ajustes realizados al esquema de base de datos 
	conforme evolucionaron los requerimientos del sistema.
	
	\item Implementar el módulo de Usuarios con sus métodos personalizados de creación 
	y actualización vinculados a roles, y con guards de acceso que diferencien usuarios 
	activos de administradores.
	
	\item Implementar el módulo de Autenticación, centralizando la emisión y validación de tokens 
	JWT para todas las rutas protegidas del sistema.
	
	\item Implementar el módulo de Catálogo para las 
	entidades de referencia del sistema: \texttt{roles}, \texttt{brands}, 
	\texttt{models}, \texttt{categories}, \texttt{statuses} y \texttt{ubications}.
	
	\item Implementar el módulo de Productos con sus campos de relación opcionales como \texttt{brand\_id}, 
	\texttt{model\_id} y \texttt{category\_id}.
	
	\item Implementar el módulo de Dispositivos con métodos personalizados de consulta, 
	filtrado y cambio de estado, incorporando un registro automático de historial.
	
	\item Implementar el módulo de Préstamos con el flujo completo de ciclo de vida 
	de una solicitud: creación, aprobación o rechazo, entrega y devolución, 
	actualizando los estados de los dispositivos involucrados.
\end{itemize}

% -------------------------------------------------------
% HITO 4: Implementación del frontend
% -------------------------------------------------------

\subsection*{Hito 4 --- Objetivo general}

Implementar el frontend completo del sistema mediante un plan de seis fases 
progresivas, desarrollando las vistas, componentes, stores y servicios necesarios 
para que administradores y usuarios puedan interactuar con todos los módulos del 
sistema a través de una interfaz funcional, coherente y adaptada a distintos 
dispositivos y resoluciones de pantalla.

\subsection*{Hito 4 --- Objetivos específicos}

\begin{itemize}
	\item Implementar los layouts, el flujo 
	de autenticación con JWT, y el control de acceso basado en roles mediante guards 
	de navegación.
	
	\item Desarrollar las vistas para 
	la administración de usuarios y roles.
	
	\item Implementar la gestión completa del catálogo de 
	referencia.
	
	\item Desarrollar las vistas del módulo de productos con distinción de roles y navegación integrada hacia 
	el módulo de dispositivos.
	
	\item Implementar las vistas del módulo de dispositivos, habilitando la gestión administrativa y la 
	interacción del usuario común con el catálogo de dispositivos disponibles.
	
	\item Implementar el ciclo de vida completo de los préstamos, conectando los endpoints de aprobación, entrega y devolución.
	
	\item Realizar una revisión de UX y UI que cubra ajustes visuales, 
	soporte de modo claro y oscuro, comportamiento \textit{responsive} en distintos 
	dispositivos y resoluciones, y retroalimentación al usuario mediante 
	notificaciones \texttt{Toast}.
\end{itemize}